% Section 3 -- Metatheory: T1-T4
% Target: 4 pages.
% Goal: state the four headline theorems and report current proof status.

DLC's metatheory is organized around four headline theorems.  This
section states each theorem in canonical form, reports its current
mechanization status, and points to the Lean proof file containing
the partial or complete proof.

\subsection{T1 -- Decidability of proof-checking}

\begin{theorem}[T1, propositional fragment]
For every context $\Gamma$, term $M$, and proposition $\varphi$ in the
propositional fragment of DLC, the question ``does $M$ prove $\varphi$
in $\Gamma$?'' is decidable.  The decision procedure runs in time
$O(|M| \cdot \log |\Gamma|)$.
\end{theorem}

\paragraph{Status: \texttt{stated}; Rust \texttt{decide\_pure} runs on the full calculus.}
A Lean function \texttt{decideLean} mirroring the Rust implementation
is forthcoming; the soundness direction (\texttt{decideLean} returns
\texttt{true} implies $\mathsf{Nonempty}(\mathsf{Deriv}~\Gamma~M~\varphi)$) is
the next mechanization target.  Completeness in Lean is deferred to
collaborator engagement.

The Rust function (\texttt{crates/dlc-core/src/decide.rs}) handles all
21 term constructors, including the no-chain-splicing condition on
\texttt{delegate}.  Nine unit tests, including a negative
\texttt{delegate\_rejects\_chain\_splicing} case, document the
implementation.

\subsection{T2 -- Cryptographic correspondence}

\begin{theorem}[T2]
For every keyring $K$ extractable from $\Gamma$:
\[
\Gamma \vdash M : \varphi \iff \Gamma \vdash_K M : \varphi
\]
That is, the logical and cryptographic typing judgments coincide.
\end{theorem}

\paragraph{Status: \texttt{stated}; conditional form planned via VCVio.}
T2 has a symbolic half and a computational half.  The symbolic half
states $\Gamma \vdash M : \varphi$ implies $\Gamma \vdash_K M : \varphi$
when all signatures verify, and conversely; this is structural
induction.  The computational half states that an adversary
distinguishing the two judgments would break Ed25519's EUF-CMA
security; this is a game-hopping reduction.

The companion artifact's \texttt{models/easycrypt/Game.eca} states the
reduction skeleton (EUF-CMA game, symbolic non-forgery game, four-step
game-hop sequence sketched).  An alternative path---closing T2
entirely within Lean using VCVio (Tuma et al.,
2026)~\cite{tuma2026vcvio}---is the autonomous-arc target.  Either
route yields the same theorem; VCVio is preferred for staying within
a single proof assistant.

\subsection{T3 -- Non-interference under delegation}

\begin{theorem}[T3]
For every IFC label $\ell_{\mathit{low}}$ and every well-typed
derivation
$\Gamma \vdash M : \varphi$ in which all hypotheses of $\Gamma$ sit at
labels $\geq \ell_{\mathit{low}}$, the term $M$ does not reveal information
to an observer at $\ell_{\mathit{low}}$.  Formally, $M$ is related to
itself by the logical relation
$\mathsf{Indistinguishable}_{\ell_{\mathit{low}}}~\varphi$.
\end{theorem}

\paragraph{Status: \texttt{stated}; atomic-fragment closure planned.}
T3 follows the standard logical-relations recipe: define
$\mathsf{Indistinguishable}_{\ell_{\mathit{low}}}~\varphi~M~M'$ by induction
on $\varphi$; prove the fundamental lemma by induction on
$\mathsf{Deriv}$; conclude non-interference as a corollary.

iris-tini (Gregersen \& Bay, POPL
2021)~\cite{gregersen2021tini} provides the reference proof technique
in Coq+Iris.  DLC's setting is simpler (no higher-order state, no
recursive types), so a plain Lean construction without Iris suffices.
This is autonomously reachable; cf. our companion document
\texttt{spec/t2-t3-lean-feasibility-2026-05.md} for the analysis.

\subsection{T4 -- Obligation soundness}

\begin{theorem}[T4, non-introduction direction]
For every reduction $M \rhd M'$ and every obligation $o$:
\[
o \in \mathsf{pendingObligations}(M') \implies o \in \mathsf{pendingObligations}(M)
\]
That is, reduction never introduces an obligation that was not
syntactically present beforehand.
\end{theorem}

\paragraph{Status: \texttt{proven\_partial}; \texttt{pendingObligations\_shift} closed.}
The supporting lemma $\mathsf{pendingObligations}~(\mathsf{shift}~t~\delta~c)
= \mathsf{pendingObligations}~t$ is fully proven by 22-case induction
on $t$ (file: \texttt{lean/DLC/ObligationSoundness.lean}).  The
substitution sub-lemma and the case analysis over $\mathsf{step}$ are the
next mechanization steps.  Prior closure attempt at PR \#16 commit
\texttt{b2768fb} (340 LOC) failed on tactic-detail bugs that smaller
PRs would have caught; the current strategy splits the proof into a
substitution-subset lemma PR and a step-case-analysis PR.

The full obligation-soundness statement (tracking multiset
multiplicities and discharge consumption) is deferred to the M3 runtime
layer where \texttt{Discharged} evidence is concrete.

\subsection{Summary of mechanization status}

\begin{table}[h]
\centering
\begin{tabular}{lll}
\toprule
Theorem & Status & File \\
\midrule
T1 (propositional, soundness) & \texttt{stated} (Rust impl. runs) & \texttt{Decidability.lean} \\
T2 (full) & \texttt{stub} & \texttt{Correspondence.lean} \\
T2 (symbolic half) & implicit in \texttt{Tamarin+ProVerif} & \\
T3 (full) & \texttt{stated} & \texttt{NonInterference.lean} \\
T4 (non-introduction) & \texttt{proven\_partial} & \texttt{ObligationSoundness.lean} \\
Substitution composition & \texttt{stated}; 4 structural lemmas proven & \texttt{Subst.lean} \\
Subject reduction & \texttt{stated} & \texttt{Reduce.lean} \\
Graded comonad laws & \texttt{proven} (3 laws) & \texttt{Graded.lean} \\
Indexed monad laws & \texttt{proven} (4 laws incl.\ strength) & \texttt{IndexedMonad.lean} \\
\bottomrule
\end{tabular}
\caption{Mechanization status as of May 2026.  Twelve theorems are
fully proven with no \texttt{sorry} and no new axioms.  Four headline
theorems are stated with proof closures detailed in
Section~\ref{sec:conclusion}.}
\label{tab:status}
\end{table}
